
%\needspace{4\baselineskip}
\color{uniblue}{\subsection{Фиксация положения выключателей}}
\color{black}

\begin{enumerate}[label=\arabic{section}.\arabic{subsection}.\arabic{enumi}, labelsep=4pt, leftmargin=0pt, itemindent=57pt, itemsep=0pt, parsep=0pt] % если что можно крутить itemsep

\item
Функциональный блок фиксации положения выключателей (ФПВ) предназначен для определения состояния выключателей на основании технологической информации:

\begin{itemize}
    \item состояния блок-контактов выключателя;
    \item положения разъединителей в цепи выключателя;
    \item информации о выводе в ремонт выключателя.
\end{itemize}

\item
Контроль положения разъединителей при формировании сигналов включенного или отключенного положения выключателей (например, \textit{«В1 включен»} и \textit{«В2 отключен»} для выключателя В1) вводится программным ключом \textbf{«Контр. Р В1»}. Также имеется возможность выбрать количество разъединителей, участвующих в формировании сигналов включенного или отключенного положения выключателей, для выключателя В1 она задается программным переключателем \textbf{«Схема Р В1»}.

\item
Также функция ФПВ пофазно контролирует состояние блок-контактов выключателя и в случае расхождения полюсов формирует сигнал \textit{«Внутр. пуск. ЗНФ В1»}.

\item
Функциональная схема логики ФПВ В1 приведена на рисунке \ref{pict_fpv:fpv}. Функциональная схема логики ФПВ В2 выполняется аналогично схеме логики ФПВ В1.

\vspace{2mm}
\begin{figure}[H]
\centering
\includegraphics[width=0.9\textwidth,height=0.9\textheight,keepaspectratio]{\fbpath/761.3381 ФПВ/img/fpv.pdf}
\captionof{figure}{Функциональная схема логики ФПВ В1}\label{pict_fpv:fpv}
\end{figure}

\item
Уставки ФПВ В1 представлены в таблице \ref{desc_fpv:tbl1}. Уставки ФПВ В2 аналогичны уставкам ФПВ В1.

\small
\begin{longtable}{|>{\centering\arraybackslash}m{5.3cm}|>{\centering\arraybackslash}m{3.3cm}|>{\centering\arraybackslash}m{4.2cm}|>{\centering\arraybackslash}m{1.8cm}|>{\centering\arraybackslash}m{1cm}|}
\caption{Параметры для настройки ФПВ В1\hfill\vspace{-0.5\baselineskip}}\label{desc_fpv:tbl1}\\ 
\hline
\rowcolor{gray!30}
Параметр & Условное обозначение на схеме & Значение/ Диапазон & Единица измерения & Шаг \\ 
\hline
\endfirsthead
\caption*{\hspace{3pt}\emph{Продолжение таблицы \ref{desc_fpv:tbl1}\hfill\vspace{-0.5\baselineskip}}} \\ % сделано по ГОСТ 2.105 п.6.8.7
\hline
\rowcolor{gray!30}
Параметр & Условное обозначение на схеме & Значение/ Диапазон & Единица измерения & Шаг \\ 
%\hline
\endhead

%===m>ФПВ В1|ФПВ
\raggedright \textcolor{red}{FullDescription НЕ НАЙДЕН!} & \centering Контр. Р В & \centering Введено \\ Выведено & \centering --- & \centering \arraybackslash --- \\
\hline
%===m
\end{longtable}
\normalsize

\end{enumerate}